% Source-aligned English for printed p.202.
\clearpage
\SourcePage{202}

\LayerTarget{tohoku:proof:proposition-5.2.2:p1}
\noindent functor from ${\mathbf C}^{X(G)}$ to ${\mathbf C}^Y$. But, by the computation of these
derived functors in Lemma 3.7.2, and recalling that an injective object of
${\mathbf C}^{X(G)}$ is a flasque sheaf, one sees that the same result is obtained as
when $f_*$ is regarded as a functor from ${\mathbf C}^X$ to ${\mathbf C}^Y$, i.e.
\textit{$R^qf_*(A)$ is the sheaf on $Y$ associated with the presheaf formed
by the $H^q(f^{-1}(U),A)$.}

\LayerTarget{tohoku:remark:cohomological-hypotheses-5.2}
The hypothesis of Proposition 5.2.2 must be regarded as a cohomological
equivalent of the hypothesis that $G$ is a ``discontinuous group of
operators'' on $X$. Another important case is that in which
$\mathbf H^q(G,A)=0$ for $q>0$, a hypothesis that must be regarded as a
cohomological equivalent of the usual condition that ``$G$ acts without
fixed points''.

\LayerTarget{tohoku:proposition:5.2.3}
\textsc{Proposition 5.2.3.} \textit{Suppose that
$\mathbf H^q(G,A)=0$ for $q>0$. Then $H^n(X;G,A)$ is isomorphic to
$H^n(Y,A^G)$; hence $H^*(Y,A^G)$ is the abutment of a cohomological spectral
sequence whose initial term is
$II_2^{p,q}(A)=H^p(G,H^q(X,A))$.}

\LayerTarget{tohoku:proof:proposition-5.2.3}
This is an immediate consequence of the first spectral sequence of Theorem
5.2.1.
\LayerTarget{tohoku:application:classical-spaces-with-operators}
We thus obtain the classical form of the theory of spaces with operators [4].
In the next section we shall give conditions under which this proposition
holds, where $G$ is a ``discontinuous group of operators on $X$ without fixed
points''. Other such conditions are related to the characteristic of a ground
field $k$ for the sheaf $A$ and the orders of the stabilizers $G_y$. A
particularly simple case is the following (which moreover is proved directly
and very easily, and is probably well known):

\LayerTarget{tohoku:corollary:5.2.3}
\textsc{Corollary.} \textit{Suppose, in addition to
$\mathbf H^q(G,A)=0$ for $q>0$, that $G$ is a finite group of order $m$, that
$X$ is separated, and that multiplication by $m$ on $A$ is an automorphism of
$A$ (for example, $A$ is a sheaf of vector spaces over a field whose
characteristic does not divide $m$). Then
$H^n(Y,A^G)=H^n(X,A)^G$, both sides being isomorphic to $H^n(X;G,A)$.}

\LayerTarget{tohoku:transition:final-particular-case-5.2}
One final particular case of interest is the following.

\LayerTarget{tohoku:proposition:5.2.4}
\textsc{Proposition 5.2.4.} \textit{Suppose that $H^q(X;A)=0$ for
$q>0$. Then $H^n(X;G,A)=H^n(G,\Gamma(A))$; consequently
$H^*(G,\Gamma(A))$ is the abutment of a spectral sequence whose initial term
is $I_2^{p,q}(A)=H^p(Y,\mathbf H^q(G,A))$.}

\LayerTarget{tohoku:application:modular-group-cohomology}
This proposition can be used to compute the cohomology of certain groups, for
example the constant-coefficient cohomology of various modular groups in one
variable [10].

\LayerTarget{tohoku:remark:5.2.4:1}
\textsc{Remarks. 1.} Suppose that a $G$-sheaf ${\mathbf O}$ is given and that
we restrict attention to $G$-${\mathbf O}$-modules, $\Gamma_X^G$ then being regarded as
a functor from ${\mathbf C}^{\mathrm O(G)}$ to the category of modules over $\Gamma({\mathbf O})^G$.
Its derived functors coincide with the preceding functors $H^p(X;G,A)$, as
follows at once from the fact that every injective object of ${\mathbf C}^{\mathrm O(G)}$ is
annihilated by the $H^n(X;G,A)$ for $n>0$. This last fact will be proved in
Lemma 5.6.1 below (where $B$ is replaced by $A$, and $A$ by ${\mathbf O}$).

\LayerTarget{tohoku:remark:5.2.4:2}
2. From the two exact sequences (5.2.8) and (5.2.9), one can
eliminate $H^*(X;G,A)$ so as to obtain relations between the cohomology of
$Y$, that of $X$

\TohokuContinuation{203}
